\documentclass[a4paper,11pt]{article}

% Use UTF-8 for input encoding
\usepackage[utf8]{inputenc}

% Use T1 font encoding
\usepackage[T1]{fontenc}

% Hyphenation for supported languages
\usepackage[english,ngerman]{babel}

% Floats
\usepackage{float}

% Advanced tables
\usepackage{array}

% Mathematics
\usepackage{amsmath}
\usepackage{amssymb}

% TikZ pictures
\usepackage{tikz}

% Source code listing
\usepackage{listings}
\lstset{
basicstyle=\ttfamily,
frame=single,
captionpos=b
}

% Conditional branching
\usepackage{ifthen}

% Multilanguage support
% (see http://en.wikibooks.org/wiki/LaTeX/Internationalization#Multilingual_versions)
\def\addlocale#1#2{
\ifthenelse{ \equal{\locale}{#1} }{
  \selectlanguage{#2}
  \expandafter\long\expandafter\def\csname#1\endcsname ##1{##1}
  }{
  \expandafter\long\expandafter\def\csname#1\endcsname ##1{}
  }
}

\AtBeginDocument{%
  % Select document language
  \en{\selectlanguage{english}}%
  \de{\selectlanguage{ngerman}}%
}%

% Define supported languages
\addlocale{en}{english}
\addlocale{de}{ngerman}

% Support for hyperlinks
\usepackage[
  hidelinks,
  pdftitle={\en{ERCIS beamer style}\de{ERCIS Beamervorlage}},
  pdfsubject={},
  pdfauthor={Dominik Lekse},
  pdfkeywords={}
]{hyperref}

% Styles
% Use smaller margins (set the margins to one inch)
\usepackage[margin=2.54cm]{geometry}

% Indention
\setlength{\parindent}{0cm}

% List formatting
\renewcommand{\labelitemi}{\raisebox{0.75pt}{\scriptsize$\blacksquare$}}

% Define ERCIS colors (see official corporate design manual)
\definecolor{ercisblack}{RGB}{0,0,0}
\definecolor{ercisgrey}{RGB}{94,94,93}
\definecolor{ercisred}{RGB}{133,35,57}
\definecolor{ercislightred}{RGB}{200,156,166}
\definecolor{ercisblue}{RGB}{135,151,163}
\definecolor{ercisclaim}{named}{ercislightred}

% Additional ERCIS colors
\definecolor{ercisdarkblue}{RGB}{67,92,139}
\definecolor{erciscyan}{RGB}{0,156,179}
\definecolor{ercisorange}{RGB}{231,124,18}
\definecolor{ercisgreen}{RGB}{135,191,42}

% Define WWU colors (see official Corporate Design manual, p. 21)
\definecolor{pantoneblack7}{RGB}{62,62,59}


% Document variables
\newcommand{\varstylename}{ercisbeamer}
\newcommand{\varstylefile}{\varstylename.sty}
\newcommand{\vartemplatefile}{template\textbackslash presentation.tex}

% Meta information
\title{\en{ERCIS beamer template}\de{ERCIS Beamervorlage}}
\author{Dominik Lekse\\
lekse@lugmedia.de}
\date{\en{Last update}\de{Letzte Aktualisierung}: 10.06.2013}

\begin{document}
% ERCIS Logo
\begin{center}
  \includegraphics[height=4cm]{images/ercis-logo.pdf}
\end{center}

% Title
% \maketitle
{\let\newpage\relax\maketitle}

% Table of contents
\renewcommand{\contentsname}{\en{Contents}\de{Inhalt}}
\tableofcontents
\pagebreak

\section{\en{Prerequisites}\de{Voraussetzungen}}
\en{Prior to the use of the style, the following requirements must be satisfied to ensure proper rendering of presentations.}
\de{Für eine einwandfreie Darstellung der Präsentationen müssen vor der Verwendung der Vorlage die folgenden Voraussetzungen erfüllt sein.}

\begin{itemize}
  \item
    \en{Install the typeface \emph{FF Meta}, which is available for members of the University of Münster in the \href{https://www.uni-muenster.de/marketing/corporatedesign/schriften.html}{corporate design manual}.}
    \de{Installieren Sie die Schriftart \emph{FF Meta}, die Angehörigen der Universität Münster im \href{https://www.uni-muenster.de/marketing/corporatedesign/schriften.html}{Corporate Design Manual} zur Verfügung steht.}
\end{itemize}

\section{\en{Usage}\de{Verwendung}}
\en{To create a presentation using the template, the following issues have to be considered.}
\de{Bei der Erstellung einer neuen Präsentation mit der Vorlage, müssen die folgenden Hinweise beachtet werden.}

\begin{itemize}
  \item
    \en{Copy the style \emph{\varstylefile} into the folder which contains the source files of your presentation.}
    \de{Kopieren Sie den Stil \emph{\varstylefile} in den Ordner, der die Quelldateien ihrer Präsentation enthält.}
  \item
    \en{Add the line \texttt{\textbackslash usepackage[...]\{\varstylename\}} including the appropriate options to the preamble of your presentation.}
    \de{Fügen Sie \texttt{\textbackslash usepackage[...]\{\varstylename\}} mit den entsprechenden Optionen der Präambel Ihrer Präsentation hinzu.}
  \item
    \en{Create the sub folders \emph{beamer-cache}, \emph{handout-cache} and \emph{article-cache} inside your presentation folder. These folders are used by the template to cache embedded graphics depending on the current mode (\emph{beamer}, \emph{handout} or \emph{article}). For more information on the caching mechanism, refer to section~\ref{sec:features-caching}.}
    \de{Erstellen Sie im Präsentationordner die Unterordner \emph{beamer-cache}, \emph{handout-cache} sowie \emph{article-cache}. In diesen Ordnern werden abhänig vom aktuellen Modus (\emph{Beamer}, \emph{Handout} oder \emph{Article}) für das Präsentationlayout benötigte bzw. Ihre eigenen TikZ-Grafiken zwischengespeichert. Weitere Informationen zum Caching-Mechanismus können Sie Abschnitt~\ref{sec:features-caching} entnehmen.}
  \item
    \en{Compile the presentation using the command line argument \emph{-shell-escape}. This is required by the \LaTeX\ compiler to invoke sub processes for caching individual graphics.}
    \de{Kompilieren Sie Ihre Präsentation mit dem Argument \emph{-shell-escape}. Dies erlaubt dem \LaTeX\ Kompiler das Starten von Unterprozessen für die Zwischenspeicherung einzelner Grafiken.}
\end{itemize}

\en{In addition to the style, the distribution includes a template \emph{\vartemplatefile}, which can be used as a starting point when creating a new presentation.}
\de{Das Vorlagenpaket beinhaltet neben dem Stil auch eine Vorlage \emph{\vartemplatefile}, die für die Erstellung einer neuen Präsentation verwendet werden kann.}

\section{\en{Features}\de{Funktionen}}

\en{An overview of the features provides by the template is given in section~\ref{sec:features-overview}. In addition, the distribution of the template comprises the source \emph{presentation\textbackslash presentation.tex}, which contains the use of elements commonly used in presentation.}
\de{Eine Übersicht der Funktionen der Vorlage ist in Abschnitt~\ref{sec:features-overview} dargestellt. Zusätzlich befindet sich im Vorlagenpaket das Dokument \emph{presentation\textbackslash presentation.tex}, im welchem exemplarisch die Verwendung von häufig benötigten Präsentationselementen veranschaulicht.}

\subsection{\en{Overview}\de{Übersicht}}
\label{sec:features-overview}

\begin{itemize}
  \item
    \en{Support for all elements of the beamer package}
    \de{Unterstützung aller Gestaltungselemente des Beamer-Pakets}
  \item
    \en{Multiple display formats}
    \de{Mehrere Anzeigeformate}
  \item
    \en{Caching of TikZ pictures}
    \de{Zwischenspeichern von TikZ Grafiken}
\end{itemize}

\subsection{\en{Display formats}\de{Anzeigeformate}}
\label{sec:features-formats}

\en{The template supports three common display formats, which are summarized in the following table~\ref{tab:features-formats}. To define the display format of a presentation, the corresponding option have to be added to the option list of the \texttt{\textbackslash usepackage[...]\{\varstylename\}} command in its preamble.}
\de{Die Vorlage unterstützt drei verschiedene Anzeigeformate, die in der folgenden Tabelle~\ref{tab:features-formats} zusammengefasst sind. Um das Anzeigeformat einer Präsentation festzulegen, muss die entsprechende Option zu der Liste der Optionen innerhalb Befehls \texttt{\textbackslash usepackage[...]\{\varstylename\}} hinzugefügt werden. Falls kein Format explizit angegeben wird, werden Präsentationen im 16:9-Format erstellt.}

\tikzstyle{ratiobox}=[draw, rectangle, text centered, inner sep=0mm, outer sep= 0mm]

\begin{table}[H]%
    \begin{center}%
        \renewcommand{\arraystretch}{1.25}
        \begin{tabular}{>{\centering} m{2cm} >{\centering} m{2cm} >{\centering} m{2cm} >{\centering} m{2cm}}
            \hline
            \textbf{\en{Aspect ratio}\de{Format}} & \textbf{\en{Option}\de{Option}} & \textbf{\en{Width}\de{Breite}} & \textbf{\en{Height}\de{Höhe}} \tabularnewline
            \hline\noalign{\vskip 6pt}
            \tikz{\node [ratiobox, text width=15mm, minimum height=11.25mm] {4:3};} & \texttt{small43} & 128 mm & 96 mm \tabularnewline
            \tikz{\node [ratiobox, text width=15mm, minimum height=9.375mm] {16:10};} & \texttt{wide10} & 144 mm & 90 mm \tabularnewline
            \tikz{\node [ratiobox, text width=15mm, minimum height=8.4375mm] {16:9};} & - & 160 mm & 90 mm \tabularnewline
            \hline
        \end{tabular}
        \caption{\en{Supported display formats}\de{Unterstützte Anzeigeformate}\label{tab:features-formats}}%
        \renewcommand{\arraystretch}{1}
    \end{center}%
\end{table}%

\subsection{\en{Frame layouts}\de{Folienlayouts}}


\en{The template extends the \texttt{frame} environment of the beamer package with additional options to influence the layout of a frame. These options are summarized in table~\ref{tab:features-framelayouts} can be combined with each other.}
\de{Die Vorlage erweitert die \texttt{frame} Umgebung des beamer Pakets um zusätzliche Optionen, die das Layout einer Folie beeinflussen. Diese zusätzlich verfügbaren Optionen sind in Tabelle~\ref{tab:features-framelayouts} zusammengefasst und können beliebig miteinander kombiniert werden.}

\en{Optionally, frames will display a single line containing a (hyperlinked) list of section titles at the very top. This is enabled by passing the \textt{navigation} option to the \varstylename package (e.\,g., \texttt{\textbackslash usepackage[navigation]\{\varstylename\}}. This line can then be suppressed for individual frames by adding the \textt{nonav} option (see table~\ref{tab:features-framelayouts}).}
\de{Inhaltsfolien können optional um eine einzeilige Abschnitts-Navigation ergänzt werden. Diese wird angezeigt, indem dem \varstylename-Paket die \textt{navigation}-Option übergeben wird (z.\,B. \texttt{\textbackslash usepackage[navigation]\{\varstylename\}}). Diese Navigation kann dann für einzelne Folien wieder unterdrückt werden, indem ihnen die \textt{nonav}-Option gegeben wird (s. Tabelle~\ref{tab:features-framelayouts}).}

\begin{table}[H]%
  \begin{center}%
    \renewcommand{\arraystretch}{1.25}
    \begin{tabular}{>{\centering} m{3cm} m{8cm}}
      \hline
      \textbf{\en{Option}\de{Option}} & \textbf{\en{Layout effect}\de{Layouteffekt}} \tabularnewline
      \hline
      \texttt{nologo} & \en{Hides the ERCIS logo in the headline}\de{Unterdrückt das ERCIS Logo in der Kopfzeile} \tabularnewline
      \texttt{notitle} & \en{Hides the title in the headline}\de{Unterdrückt den Titel in der Kopfzeile} \tabularnewline
      \texttt{nosubtitle} & \en{Hides the subtitle in the headline}\de{Unterdrückt den Untertitel in der Kopfzeile} \tabularnewline
      \texttt{nofootline} & \en{Hides the footline}\de{Unterdrückt die Fußzeile} \tabularnewline
      \texttt{nonav} & \en{Hides the section navigation}\de{Unterdrückt das Inhaltsverzeichnis} \tabularnewline
      \hline
    \end{tabular}
    \caption{\en{Additional frame layouts}\de{Zusätzliche Folienlayouts}\label{tab:features-framelayouts}}%
    \renewcommand{\arraystretch}{1}
  \end{center}%
\end{table}%

\subsection{\en{Section Frame}\de{Abschnittsfolie}}

\en{The template offers a standard template for section frames. An example for this slide is seen here:}
\de{Die Vorlage bietet eine Standardvorlage für Abschnittsfolien an. Ein Beispiel dafür ist:}
\begin{lstlisting}[caption=\en{Example for section frame}\de{Beispiel für Abschnittsfolie}]
\section{Content structuring}
\begin{frame}[nologo]
    \setbeamertemplate{section image}
        {\includegraphics[scale=.4]{images/Demo.png}}
    \setbeamertemplate{section text}{The art of structuring}
    \sectionpage
\end{frame}
\end{lstlisting}

\en{The beamer template \texttt{section image} has to be used to insert a grayscale image to the section frame.
The beamer template \texttt{section text} can be used to add an optional description text.
The frame title is automatically taken from the current section title.
The command \texttt{\textbackslash sectionpage} will render the section frame.}
\de{Das Beamer-Template \texttt{section image} muss verwendet werden, um ein Bild in Grauschattierungen hinzuzufügen. Das Beamer-Template \texttt{section text} kann für einen optionalen Beschreibungstext verwendet werden.
Der Folientitel ist automatisch aus dem aktuellen Abschnittstitel übernommen.
Das Kommando \texttt{\textbackslash sectionpage} rendert die Abschnittsfolie.}

% \subsection{Font selection}
% TODO describe font fallback mechanism

\subsection{\en{Languages}\de{Sprachen}}
\en{The language of the presentation is defined using an option. \emph{English} is used by default if a language is not specified.}
\de{Die Sprache der Präsentation wird über eine Option festgelegt. Falls keine Option angegeben wird, ist \emph{English} als Sprache voreingestellt.}
\begin{table}[H]%
  \begin{center}%
    \renewcommand{\arraystretch}{1.25}
    \begin{tabular}{>{\centering} m{2cm} >{\centering} m{2cm}}
      \hline
      \textbf{\en{Language}\de{Spache}} & \textbf{\en{Option}\de{Option}} \tabularnewline
      \hline
      \en{English}\de{Englisch} & \texttt{english} \tabularnewline
      \en{German}\de{Deutsch} & \texttt{german} \tabularnewline
      \hline
    \end{tabular}
    \caption{\en{Supported languages}\de{Unterstützte Sprachen}\label{tab:features-languages}}%
    \renewcommand{\arraystretch}{1}
  \end{center}%
\end{table}%

\subsection{\en{Colors}\de{Farben}}
\en{The template provides the following named colors for use inside a presentation. These colors colors comply with the ERCIS corporate design.}
\de{Die Vorlage stellt die folgenden benannten Farben für eine Verwendung in einer Präsentation zur Verfügung. Diese Farben entsprechen dem ERCIS Corporate Design.}

\begin{table}[H]%
  \begin{center}%
    \tikzstyle{colorsample}=[rectangle, minimum width=7.5mm, minimum height=7.5mm, use as bounding box, outer sep=2em]
    \renewcommand{\arraystretch}{1.25}
    \begin{tabular}{>{\centering} m{2cm} >{\centering} m{3cm} >{\centering} m{5cm}}
      \hline
      \textbf{\en{Color}\de{Farbe}} & \textbf{\en{Name}\de{Name}} & \textbf{\en{RGB / Hexadecimal}\de{RGB / Hexadezimal}} \tabularnewline
      \hline\noalign{\vskip 6pt}
      \tikz{\node[colorsample, fill=ercisblack] {};} & ercisblack & 0 0 0 / \#000000 \tabularnewline
      \tikz{\node[colorsample, fill=ercisgrey] {};} & ercisgrey & 94 94 93 / \#5e5e5d \tabularnewline
      \hline\noalign{\vskip 6pt}
      \tikz{\node[colorsample, fill=ercisred] {};} & ercisred & 133 35 57 / \#852339 \tabularnewline
      \tikz{\node[colorsample, fill=ercislightred] {};} & ercislightred & 200 156 166 / \#c89ca6 \tabularnewline
      \tikz{\node[colorsample, fill=ercisblue] {};} & ercisblue & 135 151 163 / \#8797a3 \tabularnewline
      \hline\noalign{\vskip 6pt}
      \tikz{\node[colorsample, fill=ercisdarkblue] {};} & ercisdarkblue & 67 92 139 / \#435c8b \tabularnewline
      \tikz{\node[colorsample, fill=erciscyan] {};} & erciscyan & 0 156 179 / \#009cb3 \tabularnewline
      \tikz{\node[colorsample, fill=ercisorange] {};} & ercisorange & 231 124 18 / \#e77c12 \tabularnewline
      \tikz{\node[colorsample, fill=ercisgreen] {};} & ercisgreen & 135 191 42 / \#87bf2a \tabularnewline
      \hline\noalign{\vskip 6pt}
      \tikz{\node[colorsample, fill=pantoneblack7] {};} & pantoneblack7 & 62 62 59 / \#3e3e3b \tabularnewline
      \hline
    \end{tabular}
    \caption{\en{Supported colors}\de{Unterstützte Farben}\label{tab:features-colors}}%
    \renewcommand{\arraystretch}{1}
  \end{center}%
\end{table}%

\subsection{\en{Multi-line title}\de{Mehrzeiliger Titel}}
\label{sec:features-multilinetitle}
\en{You can use a multi-line title. If the spacing between the lines does not look as expected, insert manual \texttt{\textbackslash titlenewline} commands at the respective spots of the title.}

\de{Mehrzeilige Titel können verwendet werden. Falls der Zeilenabstand zwischen den Zeilen nicht identisch ist kann dies mit dem Befehl \texttt{\textbackslash titlenewline} angepasst werden. Dieser muss dafür manuell an den entsprechenden Stellen im Titel eingefügt werden.}

\subsection{\en{Caching of TikZ pictures}\de{Zwischenspeichern von TikZ Grafiken}}
\label{sec:features-caching}
\en{By default, all TikZ pictures embedded in a presentation are cached using the externalization mechanism of the PGF/TikZ package.\footnote{For detailed information on the externalization mechanism, refer to chapter 32, pp.~343ff. of the \href{http://www.ctan.org/tex-archive/graphics/pgf/base/doc/generic/pgf/pgfmanual.pdf}{TikZ and PGF manual}.} The main idea of this mechanism is to render a TikZ picture only once and cache it in a separate PDF document. In subsequent compilations, the cached PDF document substitutes the original TikZ picture instead of rendering it again.

The template required this mechanism since it internally uses it to generate the logos on the title and final page as well as on the headline during the first compile run.}
\de{In der Voreinstellung werden alle in einer Präsentation eingebetteten TikZ-Grafiken mittels der Externalization-Funktion des PFG/TikZ Pakets zwischengespeichert.\footnote{Für weitere Informationen zur Externalization-Funktion, siehe Kapitel 32, pp.~343ff. des \href{http://www.ctan.org/tex-archive/graphics/pgf/base/doc/generic/pgf/pgfmanual.pdf}{TikZ/PGF Handbuchs}. } Die Grundidee dieser Funktion ist, TikZ Grafiken einmalig zu rendern und in einem separaten PDF-Dokument zwischenzuspeichern. In weiteren Kompiliervorgängen wird dann das zwischenzuspeicherte PDF-Dokument eingebunden anstatt die zugehörige TikZ-Grafik erneut zu rendern. Dadurch können Präsentationen, in denen viele TikZ-Grafiken eingebettet sind, schneller erstellt werden.

Die Vorlage setzt diese Funktion unbedingt voraus, da sie intern für die Erstellung der Logos auf der Titel- und Schlussseite sowie in der Kopfzeile verwendet wird.}

\begin{itemize}
  \item
    \en{Prepend a picture with \texttt{\textbackslash tikzsetfigurename\{name-of-picture-\}} to name the cached file accordingly.}
    \de{Verwenden Sie den Befehl \texttt{\textbackslash tikzsetfigurename\{name-of-picture-\}} um das zwischengespeicherte PDF-Dokument zu benennen.}
  \item
    \en{TikZ pictures which contain the \texttt{overlay} option cannot be cached and must be prepended with \texttt{\textbackslash nocachePicture}. Try to avoid overlay pictures since they require two compile runs to be positioned properly.}
    \de{Eine Zwischenspeicherung von TikZ Grafiken mit der \texttt{overlay} Option ist nicht möglich und muss durch das Voranstellen von \texttt{\textbackslash nocachePicture} unterdrückt werden. Versuchen Sie Overlay-TikZ-Grafiken zu vermeiden, da sie außerdem zwei Komplierungsdurchgänge benötigen um korrekt positioniert zu werden.}
\end{itemize}

\en{To support the iterative workflow of creating TikZ pictures in a presentation, the template provides some macros as shortcuts for influencing the caching of a particular picture. These macros are summarizes in table~\ref{tab:features-tikz-caching-macros} and can be prepended to a TikZ picture in order to achieve the corresponding effect. For example, prepend \texttt{\textbackslash nocachePicture} to supress the caching of a single picture while creating and testing it. Once it has been finished, remove or put the macro in comments to cache the final version of the picture.}
\de{Die Vorlage stellt zur Untersützung des iterativen Erstellungsprozesses von TikZ Grafiken einige Makros als Kurzbefehle bereit. Diese Makro sind in der folgenden Tabelle~\ref{tab:features-tikz-caching-macros} zusammengefasst und können einer TikZ Grafik vorangestellt werden um den entsprechenden Effekt zu erreichen. Beispielsweise kann während der Erstellung einer TikZ Grafik das Makro \texttt{\textbackslash nocachePicture} vorangestellt werden, so dass das Zwischenspeichern der Grafik unterdrückt wird. Sobald die TikZ Grafik fertiggestellt ist, kann das Makro entfernt bzw. auskommentiert werden.}

\begin{table}[H]%
  \begin{center}%
    \renewcommand{\arraystretch}{1.25}
    \begin{tabular}{m{3.5cm} m{6cm} m{5.5cm}}
      \hline
      \textbf{\en{Command}\de{Befehl}} & \textbf{\en{Name}\de{Name}} & \textbf{\en{Description}\de{Beschreibung}} \tabularnewline
      \hline
      \texttt{\textbackslash rerenderPicture} & \texttt{\textbackslash tikzset\{external/remake next=true\}} & \en{Force the rerendering of the subsequent picture}\de{Erzwingt das erneute Rendern und Zwischenspeichern der nachfolgenden Grafik} \tabularnewline
      \texttt{\textbackslash rerenderAll} & \texttt{\textbackslash tikzset\{external/force next=true\}} & \en{Force the rerendering of all pictures}\de{Erzwingt das erneute Rendern und Zwischenspeichern aller Grafiken} \tabularnewline
      \hline
      \texttt{\textbackslash cachePicture} & \texttt{\textbackslash tikzset\{external/export next=true\}} & \en{Explicitly enable caching of the subsequent picture}\de{Aktiviert explizit das Zwischenspeichern der nachfolgenden Grafik} \tabularnewline
      \hline
      \texttt{\textbackslash nocachePicture} & \texttt{\textbackslash tikzset\{external/export next=false\}} & \en{Explicitly disable caching of the subsequent picture}\de{Verhindert das Zwischenspeichern der nachfolgenden Grafik} \tabularnewline
      \texttt{\textbackslash nocacheAll} & \texttt{\textbackslash tikzset\{external/export=false\}} & \en{Explicitly disable caching of all pictures}\de{Verhindert das Zwischenspeichern aller Grafiken} \tabularnewline
      \hline
    \end{tabular}
    \caption{\en{TikZ caching helper macros}\de{Unterstützte Farben}\label{tab:features-tikz-caching-macros}}%
    \renewcommand{\arraystretch}{1}
  \end{center}%
\end{table}%

\section{\en{Display issues}\de{Darstellungsprobleme}}

\subsection{\en{ERCIS logo in Adobe Reader and Adobe Acrobat Professional}\de{ERCIS logo in Adobe Reader und Adobe Acrobat Professional}}
\en{When viewing a presentation in Adobe Reader or Adobe Acrobat Professional, the space between the grey dots in the ERCIS logo seem to disappear. This issue is caused by the gap between these dots being too narrow in relation to the size of the logo. In Adobe Reader and Adobe Acrobat Professional, a feature to enhance such thin lines is enabled by default, but unfortunately causes the logo not to be rendered properly. To disable this feature, click \emph{Edit} > \emph{Preferences} > \emph{Page Display} and remove the checkmark from the option \emph{Enhance thin lines}.}
\de{Bei der Anzeige einer Präsentation im in Adobe Reader oder im Adobe Acrobat Professional scheinen die Zwischenräume zwischen den grauen Punkten im ERCIS Logo zu verschwinden. Dies wird durch einen im Verhältnis zur Größe des Logos zu schmalen Zwischenraum verursacht. Der Adobe Reader bzw. Adobe Acrobat Professional korrigieren in der Standardeinstellung die Darstellung von dünnen Linien und stellen dadurch die das Logo verzerrt bzw. fehlerhaft dar. Die entsprechende Funktion kann unter \emph{Bearbeiten} > \emph{Voreinstellungen} > \emph{Seitenanzeige} durch Deaktivieren der Option \emph{Dünne Linien deutlicher darstellen} unterdrückt werden.}

\end{document}
